\section{A rational multistage schedule with polynomial space}\label{sec:polyschedule}
A fixed alphabet is unnecessarily restrictive. The root is constrained by
large degrees on both sides. Immediately below it, one side consists of
fresh packet tensors, so a larger cumulative taper is possible. Later, a
smaller alphabet allows a larger taper budget but gives slower size reduction.
The schedule below makes that tradeoff explicit. Variable alphabets are
already present in the round-two packet~\cite{roundtwo}; the result here is
an auditable extension for the combined asymmetric, slope-four interface.

\subsection{Parameters and exact margins}
Use $H,g,a_d$ from~\eqref{eq:gap}. Define $A_{-1}=0$ and, for $0\le d\le5$,
$u_d=\floor{(A_d-A_{d-1})\lambda}$ according to the table.
\begin{center}
\begin{tabular}{@{}rrrr@{}}
\toprule
Depth $d$ & Alphabet $r_d$ & Cumulative leading budget $A_d$ & Individual $t_d$\\
\midrule
0 & 13 & $27/2000$ & $27/2000$\\
1 & 14 & $21/800$ & $51/4000$\\
2 & 8 & $333/10000$ & $141/20000$\\
3 & 6 & $193/5000$ & $53/10000$\\
4 & 5 & $43/1000$ & $11/2500$\\
5 & 4 & $499/10000$ & $69/10000$\\
\bottomrule
\end{tabular}
\end{center}
For $d=6+j$, put
\begin{equation}\label{eq:poly-tail}
 r_d=3,\qquad t_d=\frac{131}{40000}\left(\frac34\right)^j,
 \qquad u_d=\floor{t_d\lambda}.
\end{equation}
The tail budget sums to $131/10000$, so $\sum_{d\ge0}t_d=63/1000$.
All scheduling operations are exact rational or integer arithmetic. The
provided generator returns $r_d,u_d,a_d,g$ without constructing $2^\lambda$.

\begin{lemma}[Uniform feasibility]\label{lem:poly-feasible}
For all sufficiently large $s$, this schedule satisfies
\eqref{eq:packing-conditions} and the precision assumptions at every active
call. Active paths have $O(\log\lambda)$ depth; all retained tapers tend to
infinity uniformly.
\end{lemma}
\begin{proof}
Consider hypothetical active depths through
$d_{\max}=\ceil{\log_{10/7}\lambda}+7$. The cumulative gap is
$O((\log\lambda)^3)=o(\lambda)$. Therefore
$a_d\le(A_d+o(1))\lambda$ in the six fixed stages, and
$a_d\le(63/1000+o(1))\lambda$ in the tail.

At the root the dominant term is $Q_0^2s^{9/10}$. Its power of $s$ is
$9/10+(27/1000)\log_2 13<1$, certified by the integer inequality
$13^{27}<2^{100}$. The row term is smaller. Below the root, the retained-side
power is $9/10+A_d\log_2r_d$. All required strict inequalities have exact
certificates:
\begin{align*}
 14^{21}&<2^{80},& 3\cdot\frac{333}{10000}&<\frac1{10},&
 6^{193}&<2^{500},\\
 5^{43}&<2^{100},&2\cdot\frac{499}{10000}&<\frac1{10},&
 3^{63}&<2^{100}.
\end{align*}
Each finite stage and the whole tail have a fixed positive power margin.
For the fresh side and rounding, the looser common estimates suffice:
\[
 2A_d\log_2r_d<\frac15,\qquad A_d,A_{d-1}\le\frac{63}{1000}.
\]
Thus both $Q_d^2C_d$ and $L_dQ_d^2$ are at most
$s^{263/1000+o(1)}$ below the root. Every new packet degree is at most
$s^{63/1000+o(1)}<s^{9/10}$, closing the degree invariant. Because every
active $M$ exceeds $s/\lambda$, these fixed power margins absorb all
logarithmic factors and imply~\eqref{eq:packing-conditions}.

More quantitatively, \eqref{eq:load} gives
$P/M=1+O(1/\lambda)+s^{-\varepsilon+o(1)}$ uniformly for some fixed
$\varepsilon>0$. Consequently the exact rational guard
\[
 P/M\le1+1/\ceil{\sqrt\lambda}
\]
passes eventually. For sufficiently large $\lambda$ it implies a child ratio
at most $7/10$, including the slowest alphabet-three tail. No active path
can exceed $d_{\max}$, validating the hypothetical depth range without a
circular assumption. Products of the sharper padded overhead over that range
are $1+o(1)$. Write
\begin{equation}\label{eq:mass-profile}
 v_0=1,\qquad v_{d+1}=v_d(2/r_d).
\end{equation}
Then $M_d\le(1+o(1))sv_d$ uniformly.

In the tail the smallest taper on that range is bounded below by
\[
 \lambda^{\,1-\ln(4/3)/\ln(10/7)+o(1)}.
\]
The exponent is positive because $4/3<10/7$. Floors therefore have uniform
relative error $o(1)$. The fixed stages have positive constant $t_d$.
All retained $u_d$ diverge, so~\eqref{eq:precision} and the local entropy
estimate apply uniformly.
\end{proof}

An implementation may use direct contraction for small inputs, zero tapers,
or failed exact guards. These fallbacks preserve correctness. The lemma shows
that they disappear from active calls eventually. The smallest margin in the
power comparisons is tiny; this theorem makes no assertion of a practical
input threshold for the complete contraction algorithm.

\subsection{Summing the recursion and certifying the coefficient}
\begin{proposition}\label{prop:poly-cost}
The polynomial-space part of Theorem~\ref{thm:main} holds with
\begin{equation}\label{eq:polyconstant}
 \Cnew=\sum_{d=0}^{5}\frac{k(r_d)v_d}{\sqrt{t_d}}
 +\frac{k(3)}{2730\sqrt{131/40000}}
   \frac{1}{1-4/(3\sqrt3)}.
\end{equation}
\end{proposition}
\begin{proof}
Along a recursive path, logarithms of tuple-enumeration costs add.
By~\eqref{eq:local-entropy} and Lemma~\ref{lem:poly-feasible}, the leading
coefficient is bounded by $\sum_d k(r_d)v_d/\sqrt{t_d}$.
At the start of the tail $v_6=1/2730$. Consecutive tail costs have ratio
$(2/3)/\sqrt{3/4}=4/(3\sqrt3)<1$, which sums to~\eqref{eq:polyconstant}.
Uniform $o(1)$ perturbations change the strictly convergent tail by $o(1)$.

Per-call interpolation, arithmetic, group selection, and bookkeeping use
$\poly(s)$ work on immutable handles. Input handling and original entry
queries retain the fixed polynomial dependence on $s+\mathcal L$ from the
source oracle interface. That entry-evaluation factor multiplies the total
query workload, rather than being reintroduced at every recursion level.
The $O(\log\lambda)$ levels therefore add only
$O(\lambda\log\lambda)$ to the logarithm of the bookkeeping cost,
which is negligible on the stated exponential scale.
Leaves cost $\exp(O(s/\lambda))$, initial tensor entries cost
$\exp(O(s^{9/10}))$, and newly created entries cost
$\exp(O(L\log s))$ with $L\le s^{63/1000+o(1)}$. All are
$\exp(o(s/\sqrt\lambda))$. Depth-first enumeration retains polynomial
space because each precision, label set, kernel description, and accumulator
has polynomial size and the depth is logarithmic in $\lambda$.
\end{proof}

The checker bounds every radical by integer square-root intervals and uses
rational intervals for $\pi$ and $\ln2$. It proves
\[
 63.781173756244<\Cnew<63.781173756245<64.
\]
It also verifies the exact integer power certificates above. An independent
coarser certificate, using $\pi<22/7$ and $\ln2>693/1000$, is
\[
 \Cnew<\frac{2297667966854588349649}{36002094515068500000}<64.
\]
Thus the convenient coefficient $64$ is not a conclusion drawn from a
floating-point search.

\subsection{Comparison with the stronger prior family}
For a fixed alphabet $r$ with $c=2/r$ and
$A_0=1/(20\log_2r)$, the reviewed asymmetric analysis already optimizes
root/tail allocation. After substituting slope-four modes, its boundary
expression is
\begin{equation}\label{eq:old-fixed}
 \frac{\pi\sqrt{2(1-1/r)}}{\ln2\sqrt{A_0}}
 \left(1+\frac{c}{(1-c^{2/3})^{3/2}}\right),
\end{equation}
where the active root-cap and other feasibility premises apply. At $r=16$
this is $66.192723693577\ldots$. The explicit inward-shifted schedule in the
asymmetric packet gives about $66.6051$ after the same kernel substitution.
An integer scan of the admissible fixed-alphabet comparisons finds its minimum
at $16$ over the recorded range; we do not promote that scan to an optimality
theorem over arbitrary schedules. Formula~\eqref{eq:polyconstant} is an
explicit feasible improvement over these stronger combined certificates.
